{\huge\chapter{Введение}}
\label{sec:Chapter0} \index{Chapter0}

\section{Актуальность темы}
Системы видеонаблюдения повсеместно применяются для анализа трафика, безопасности и городской логистики. При этом данные собираются множеством разнесённых камер с существенно различающимися условиями съёмки. Необходимость отслеживать один и тот же физический объект (автомобиль) при переходе между камерами приводит к задаче \emph{реидентификации} (ReID), в которой требуется сопоставить экземпляры объекта в разных ракурсах и в различное время. В реальных условиях распознавание регистрационных знаков бывает невозможно (ракурс, блики, дальность, правила приватности), поэтому особую важность приобретает ReID по \emph{внешнему виду}.

\section{Цель и задачи исследования}
\textbf{Цель} работы — исследовать и сопоставить методы реидентификации транспортных средств по внешнему виду при съёмке с разноракурсных камер наблюдения и построить воспроизводимый пайплайн для оценки качества на общепринятых наборах данных.

Для достижения цели решаются следующие \textbf{задачи}:
\begin{enumerate}
  \item Формализовать постановку Vehicle ReID и требования к данным и метрикам.
  \item Сформировать экспериментальную базу: выбрать и подготовить датасет VeRi‑776 как наиболее распространённый бенчмарк для Vehicle ReID.
  \item Реализовать и сравнить две модели (ResNet-50 и ViT-Small с BNNeck) и процедуры обучения (Cross‑Entropy с ослаблением меток, Triplet Loss, PK‑семплинг).
  \item Изучить влияние ключевых компонент: аугментации (Random Erasing/Color Jitter), re‑ranking (k‑reciprocal), стратегия подбора «якорь‑позитив‑негатив».
  \item Провести сравнение архитектур CNN (ResNet-50) и Transformer (ViT-Small) по метрикам CMC@k и mAP; обсудить устойчивость к доменному сдвигу между камерами.
\end{enumerate}

\section{Объект и предмет исследования}
\textbf{Объект} исследования — видеоданные многокамерных систем наблюдения в городских сценах. \textbf{Предмет} исследования — методы обучения эмбеддинг‑представлений для сопоставления изображений одного и того же автомобиля в условиях вариативности ракурса, освещения и качества.

\section{Научная новизна и практическая значимость}
\textbf{Новизна} состоит в системном сравнении и инженерной компоновке современных практик Vehicle ReID в едином воспроизводимом пайплайне с акцентом на отказ от номерных знаков и анализ влияния отдельных компонент (BNNeck, Triplet Loss, re‑ranking) при разнотипных условиях съёмки. \textbf{Практическая значимость} обусловлена применимостью результатов для городских ITS‑сценариев, где прямое распознавание номера нежелательно или невозможно.

\section{Методология и инструментарий}
В работе используются библиотека PyTorch и открытые репозитории (например, \texttt{fast-reid}/\texttt{torchreid}) для реализации моделей и протоколов обучения; воспроизводимость обеспечивается фиксированными сплитами и настройками. Оценка проводится по стандарту CMC/mAP; статистическая устойчивость проверяется повторяемыми запусками и усреднением результатов.

\section{Положения, выносимые на защиту}
\begin{enumerate}
  \item Пайплайн Vehicle ReID без использования номерных знаков обеспечивает воспроизводимое качество на VeRi‑776 при корректной настройке BNNeck, Triplet Loss и re‑ranking.
  \item Transformer-архитектура (ViT-Small) превосходит CNN-бэкбон (ResNet-50) по качеству идентификации (+5.1\% mAP, +1.3\% Rank-1) благодаря глобальному механизму внимания, при умеренном увеличении вычислительной сложности.
  \item Постобработка k‑reciprocal re‑ranking стабильно повышает mAP при минимальном накладном времени на инференс.
  \item Аугментации Random Erasing и цветовые искажения заметно повышают устойчивость к окклюзиям и доменному сдвигу.
\end{enumerate}

\section{Структура работы}
Работа состоит из введения, четырёх глав и заключения. В \S1\ref{sec:Chapter1} формулируется постановка задачи и обсуждается проблематика Vehicle ReID. В \S2\ref{sec:Chapter2} описываются используемые модели и методы обучения/оценивания. В \S3 приводятся результаты экспериментов и их анализ. В \S4 формулируются выводы и направления дальнейшей работы.
